\section{Introduction}
\label{sec:introduction}

How can a learning system construct, compose, and improve its own internal
computational mechanisms? Neural architectures already learn representations,
select information, implement memory, and organize reusable computation.
Attention, for example, learns input-dependent interactions through shared
projections and weighted aggregation~\cite{vaswani2017attention}. The question
is how computational capabilities can accumulate under continued learning:
which internal operations are acquired, how later tasks reuse them, and what
must be preserved when the parameters implementing those operations change.
An adequate answer should expose the executed computation and its learning
history, together with the resources required to retain it.

This problem has two coupled scales. At the implementation scale, a shared
parameter update changes representations, operations, and their invocation.
At the semantic scale, a capability is a complete input--computation--output
relation. Preserving an isolated operator is insufficient if its old inputs
are subsequently encoded differently or routed elsewhere. Conversely,
preserving all implementation details can prevent useful adaptation or require
separate capacity for each task. Dynamically expandable networks illustrate
that selective sharing, structural growth, and protection can be organized
jointly~\cite{yoon2018lifelong}. Their existence motivates a precise comparison
of learning organizations under shared information and resource conditions.

We investigate \emph{Prometheus}, an executable Block with a continuous internal
state, learned nonlinear residual primitives, a contextual caller, and
content-addressed memory. The caller reads the current state, input and memory
content to select or mix primitive updates. After those calls, learned
addresses and gates revise the memory's keys and values. A task supplies
observed inputs and a terminal label; gradients jointly train representation,
operations, invocation, memory and readout. No task-specific primitive table
or correct calling sequence is installed by the experimenter.

The architecture makes two interfaces explicit. First, symbolic embeddings
and modality-specific projections enter the same core, permitting direct tests
of a common computation on finite-state programs, text prefixes, audio and
image features. Second, soft mixture execution and hard-call execution use
the same parameters but follow distinct trajectories. Hard calls retain a
continuous state and readout. Their quality and costs require separate
measurement; replacing a mixture with its largest component is an empirical
change to the computation.

This continuous design follows measured limitations of earlier categorical
prototypes, including weak terminal credit and collapsed calling behavior.
Residual state updates and learned memory are concrete architectural responses,
whose effects can be tested against static calling, expanded state, recurrent
networks, Transformers and applicable state-space models. The component
operations are established mechanisms. The scientific target is an important
acquisition--retention--resource benefit from their organization, with clear
causal attribution. Historical monodromy-based protection remains an optional
research mechanism documented with its original assumptions and controls.

Current evidence consists of bounded development runs. A contextual sequence
experiment demonstrates learning with both execution interfaces; sampled
character prediction tests length changes and parameter-matched controls;
initial audio and image diagnostics exercise the continuous-input interface.
The results are mixed and have not yet established continuing capability
growth or a broad competitive advantage. This manuscript therefore connects
an implemented, inspectable architecture to a prospective continual-learning
evaluation. A validated Primary freeze and subsequent retraining separate
exploratory changes from the eventual empirical claims.
